% terms.tex — единый слой для терминов, глоссария и предметного указателя.

%% ── Индекс и глоссарий ────────────────────────────────────────────────────────
\makeindex[
  program=xindy,
  options=-L russian -C utf8,
  intoc,
  title={Предметный указатель}
]

\makeglossaries
\GlsSetXdyLanguage{english}
\GlsSetXdyCodePage{utf8}
\setglossarystyle{listgroup}
\renewcommand*{\glslistgroupheaderfmt}[1]{\textbf{#1}}

%% ── Вспомогательные макросы ──────────────────────────────────────────────────
% \DeclareTerm{label}{display}{sort}{margin note}{glossary description}
% Для склонений и нестандартного отображения в тексте используйте:
%   \term{label}
%   \term[кастомный текст]{label}
%   \Term{label}
% Звёздочная форма подавляет заметку на полях:
%   \term*{label}

\NewDocumentCommand{\DeclareTerm}{m m m m m}{%
  \newglossaryentry{#1}{%
    name={#2},
    text={#2},
    sort={#3},
    user1={#4},
    description={#5}
  }%
}

\newcommand{\termindex}[1]{%
  \index{\glsentrysort{#1}@\glsentryname{#1}}%
}

\newcommand{\termnote}[1]{%
  \marginnote{%
    \footnotesize
    \textbf{\glsentryname{#1}}\\
    \glsentryuseri{#1}%
  }%
}

\NewDocumentCommand{\term}{s o m}{%
  \termindex{#3}%
  \IfNoValueTF{#2}{%
    \ifglsused{#3}{\gls{#3}}{\gls{#3}\IfBooleanF{#1}{\termnote{#3}}}%
  }{%
    \ifglsused{#3}{\glslink{#3}{#2}}{\glslink{#3}{#2}\IfBooleanF{#1}{\termnote{#3}}}%
  }%
}

\NewDocumentCommand{\Term}{s o m}{%
  \termindex{#3}%
  \IfNoValueTF{#2}{%
    \ifglsused{#3}{\Gls{#3}}{\Gls{#3}\IfBooleanF{#1}{\termnote{#3}}}%
  }{%
    \ifglsused{#3}{\Glslink{#3}{#2}}{\Glslink{#3}{#2}\IfBooleanF{#1}{\termnote{#3}}}%
  }%
}

%% ── Базовые термины книги ────────────────────────────────────────────────────
\DeclareTerm
{pan}
{PAN}
{PAN}
{Primary Account Number — основной номер карты.}
{Primary Account Number — основной номер карты; главный идентификатор карты в платёжной инфраструктуре.}

\DeclareTerm
{iin}
{IIN}
{IIN}
{Issuer Identification Number — первые 6 или 8 цифр PAN.}
{Issuer Identification Number — первые 6 или 8 цифр PAN, идентифицирующие эмитента и карточную сеть.}

\DeclareTerm
{bin}
{BIN}
{BIN}
{Bank Identification Number — историческое название IIN.}
{Bank Identification Number — историческое название IIN, которое по-прежнему широко используется в индустрии.}

\DeclareTerm
{mii}
{MII}
{MII}
{Major Industry Identifier — первая цифра PAN.}
{Major Industry Identifier — первая цифра PAN, указывающая на отрасль эмитента.}

\DeclareTerm
{luhn}
{алгоритм Луна}
{Алгоритм Луна}
{Проверка контрольной цифры PAN; не является криптографией.}
{Алгоритм вычисления контрольной цифры PAN для обнаружения ошибок ввода; не является криптографической защитой.}

\DeclareTerm
{cvv}
{CVV}
{CVV}
{Card Verification Value — код верификации на магнитной полосе.}
{Card Verification Value — код верификации, записанный на магнитной полосе карты и вычисляемый эмитентом.}

\DeclareTerm
{sad}
{SAD}
{SAD}
{Sensitive Authentication Data — данные, которые нельзя хранить после авторизации.}
{Sensitive Authentication Data — чувствительные аутентификационные данные, которые PCI DSS запрещает хранить после авторизации.}

\DeclareTerm
{emv}
{EMV}
{EMV}
{Стандарт чиповых карт с динамической аутентификацией транзакций.}
{EMV — стандарт микропроцессорных платёжных карт, обеспечивающий динамическую аутентификацию транзакций.}

\DeclareTerm
{emvco}
{EMVCo}
{EMVCo}
{Организация, управляющая спецификациями EMV.}
{EMVCo — организация, которая поддерживает и развивает спецификации EMV.}

\DeclareTerm
{aid}
{AID}
{AID}
{Application Identifier — идентификатор платёжного приложения на чипе.}
{Application Identifier — уникальный идентификатор платёжного приложения на EMV-чипе.}

\DeclareTerm
{pse}
{PSE}
{PSE}
{Payment System Environment — каталог приложений на карте.}
{Payment System Environment — каталог платёжных приложений на карте для контактного интерфейса.}

\DeclareTerm
{ppse}
{PPSE}
{PPSE}
{Proximity Payment System Environment — каталог приложений для contactless.}
{Proximity Payment System Environment — каталог платёжных приложений на карте для бесконтактного интерфейса.}

\DeclareTerm
{gpo}
{GPO}
{GPO}
{Get Processing Options — команда начала EMV-транзакции.}
{Get Processing Options — команда EMV, с которой терминал запускает обработку выбранного приложения.}

\DeclareTerm
{aip}
{AIP}
{AIP}
{Application Interchange Profile — битовое поле возможностей приложения.}
{Application Interchange Profile — битовое поле, описывающее поддерживаемые функции EMV-приложения.}

\DeclareTerm
{afl}
{AFL}
{AFL}
{Application File Locator — список файлов и записей для чтения.}
{Application File Locator — список файлов и записей, которые терминал должен прочитать в ходе EMV-транзакции.}

\DeclareTerm
{sda}
{SDA}
{SDA}
{Static Data Authentication — статическая офлайн-аутентификация данных карты.}
{Static Data Authentication — базовый метод офлайн-аутентификации EMV-карты, проверяющий неизменность подписанных данных.}

\DeclareTerm
{dda}
{DDA}
{DDA}
{Dynamic Data Authentication — динамическая офлайн-аутентификация карты.}
{Dynamic Data Authentication — метод офлайн-аутентификации EMV, при котором карта формирует подпись на уникальных данных транзакции.}

\DeclareTerm
{cda}
{CDA}
{CDA}
{Combined DDA/Application Cryptogram Generation — аутентификация с криптограммой.}
{Combined DDA/Application Cryptogram Generation — EMV-метод, совмещающий аутентификацию карты с генерацией криптограммы транзакции.}

\DeclareTerm
{cvm}
{CVM}
{CVM}
{Cardholder Verification Method — способ проверки держателя карты.}
{Cardholder Verification Method — способ верификации держателя карты, например PIN, подпись или отсутствие проверки.}

\DeclareTerm
{arqc}
{ARQC}
{ARQC}
{Authorization Request Cryptogram — запрос карты на онлайн-авторизацию.}
{Authorization Request Cryptogram — криптограмма EMV, означающая, что карта просит онлайн-авторизацию у эмитента.}

\DeclareTerm
{tc}
{TC}
{TC}
{Transaction Certificate — офлайн-одобрение транзакции картой.}
{Transaction Certificate — криптограмма EMV, означающая офлайн-одобрение транзакции картой.}

\DeclareTerm
{aac}
{AAC}
{AAC}
{Application Authentication Cryptogram — офлайн-отклонение транзакции картой.}
{Application Authentication Cryptogram — криптограмма EMV, означающая офлайн-отклонение транзакции картой.}

\DeclareTerm
{tvr}
{TVR}
{TVR}
{Terminal Verification Results — результаты проверок терминала.}
{Terminal Verification Results — пятибайтовое EMV-поле, в котором терминал фиксирует результаты выполненных проверок.}

\DeclareTerm
{besp}
{БЭСП}
{БЭСП}
{Банковские электронные срочные платежи — RTGS-режим платёжной системы Банка России.}
{БЭСП — система банковских электронных срочных платежей, RTGS-режим платёжной системы Банка России для крупных и срочных межбанковских переводов.}

\DeclareTerm
{swift}
{SWIFT}
{SWIFT}
{Международная система передачи финансовых сообщений; не выполняет расчёт сама.}
{SWIFT (Society for Worldwide Interbank Financial Telecommunication) — международная система передачи финансовых сообщений между банками и другими финансовыми организациями; сама по себе не является расчётной системой.}

\DeclareTerm
{spfs}
{СПФС}
{СПФС}
{Система передачи финансовых сообщений Банка России — российский аналог SWIFT.}
{СПФС — Система передачи финансовых сообщений Банка России, российская инфраструктура обмена финансовыми сообщениями, функционально близкая к SWIFT.}

\DeclareTerm
{merchant}
{мерчант}
{мерчант}
{Торгово-сервисное предприятие, принимающее карточные платежи.}
{Мерчант — торгово-сервисное предприятие или онлайн-продавец, принимающий карточные платежи через эквайера.}
